\section{终章：一页纸的每日自救清单}

前面十章，浓缩成一张可以贴在墙上的纸。

\subsection{每天的三件事}

\begin{checklist}
  \item \textbf{晨间}：定下今天的「一件事」（恰当场）。把最重要的文件放到最显眼处，记录打开。别强求早起，晚上到点躺下就行；
  \item \textbf{午间}：低唤醒就翻找激发或躺下等想法冒出来，别硬坐；机械任务放音乐；冒出来的有价值想法，随手记或问 AI；
  \item \textbf{晚间}：允许自己在恰当场里被动切换；用「接下来干啥都行」骗过「必须」警报；不给自己定「今天只能玩两小时」这种精确配额。
\end{checklist}

\subsection{每周的一小时}

\begin{checklist}
  \item 给正在学的一个学科或一本书，建一次\term{模糊分类}——四五个大类，各起个名；
  \item 做一次\term{连接整理}：整理文件、收藏夹，把散落的东西连起来，删掉不满意的旧稿；
  \item 复盘：这周哪些环境、哪些毛点最容易激活你？记进大脑运行日志。
\end{checklist}

\subsection{遇到卡壳时，按这个顺序试}

\begin{enumerate}[leftmargin=2em,itemsep=2pt]
  \item 我是不是在\term{无聊}？——先处理无聊：放音乐、翻找激发、换个相关的东西看；
  \item 我是不是被\term{「必须」锁死}了？——换成「接下来两小时干啥都行」，实际只碰 A 相关；
  \item 我是不是在\term{低唤醒}状态？——别硬撑，躺下等想法冒出来，或用手机 / AI 唤醒；
  \item 我是不是\term{假性疲劳}？——撤掉「必须做 X」的念头，把文件放最显眼处，随便写几笔；
  \item 我是不是被\term{慢性压力}困死了？——换一次慕容水战法：此刻只打眼前这一仗。
\end{enumerate}

\subsection{永远记住的四句话}

\begin{checklist}
  \item 你不是不能专注，你是\term{主动专注太贵}——让被动注意力落在正事上；
  \item 你不是懒，你是\term{燃料不匹配}——你的启动器听兴趣、新奇、挑战、紧迫的话；
  \item 学历不是能力，是\term{脚手架}——缺了可以自己搭，别拿它审判自己；
  \item 自救不是变成别人，是\term{搞懂自己，然后顺着自己活}。
\end{checklist}

\subsection{最后的话}

这份讲义读到这里，你大概已经知道自己是哪一种了：可能是那个靠聪明和死线撑了二十年的「高功能者」，可能是那个被「我不行」困了很多年的「未被发现者」，更可能是夹在中间的、大多数。

无论哪种，终点都一样：\term{你不是坏掉的，你是不同频率的}。

从此你不再问「我为什么做不到」，而是问「我这套系统今天吃什么燃料、走哪条轨道、在哪个场地里跑得最好」。前者是审判，后者是工程。把审判换成工程，这份讲义就没有白读。

去把今天的那件「一件事」定下来吧。就从现在开始。
